
Four primary outcomes are anticipated from this methodology.
First, an inverse correlation is expected between the generalisation gap and rule comprehensibility: the SAM-regularised model should 
produce rules with fewer conditions than the baseline while maintaining comparable fidelity, consistent with the hypothesis that flatter 
minima yield more balanced feature representations \cite{Springer2024}.
\\\\
Second, the rules extracted from the SAM-regularised model are expected to show higher stability scores across different cancer subtypes 
in the GDSC dataset, reflecting the capture of genuine biological signal rather than training-specific noise.
\\\\
Third, the genes appearing in SAM-optimised rules are expected to show statistically significant enrichment in established KEGG pathways 
(p < 0.05) compared to baseline rules, providing biological grounding for the claim that regularisation improves not only statistical but 
clinical rule quality.
\\\\
Fourth, and contingent on the preceding outcomes, the study aims to demonstrate that the generalisation gap serves as a reliable 
diagnostic indicator of explanation fragility, and that SAM regularisation constitutes an effective intervention for improving the 
fidelity, comprehensibility, and biological validity of symbolically extracted rules in cancer pharmacogenomics.
